\RevisionAddBegin
\subsection{Encrypted, Halftone, and Learned Neighbors}
\RevisionAddEnd

\RevisionAddBegin
Encrypted-domain binary RDH uses a different carrier model, but it is relevant
to the treatment of side information. Ren, Lu, and Chen
\cite{Ren2019Encrypted} use prediction in encrypted binary images, and recent
encrypted-domain work includes adaptive block coding for grayscale carriers
\cite{Xiao2024AdaptiveBlock}. Halftone RDH also records dynamic embedding-state
information \cite{Yin2021HalftoneDESG}. These neighboring methods confirm
that side information and compression are established tools; the distinction
in this paper is their use as an exact serialized cost inside the active
mapping-selection rule.
\RevisionAddEnd

\RevisionAddBegin
Learning serves here as candidate-cost ranking, while the reversible encoder
remains deterministic and table based. Learned RDH and steganalysis show why
such ranking and detection analyses are relevant \cite{Li2021LearnedRDH,Boroumand2019SRNet,
Yedroudj2018Net,Zhang2019SRNet}. Because PEAK/ZERO substitution changes local
pattern histograms and transition rates, the experiments include grouped
detectors to delimit the payload region where marked images become
statistically separable from covers.
\RevisionAddEnd

\RevisionAddBegin
\subsection{Research Gap}
\RevisionAddEnd

\RevisionAddBegin
Table~\ref{tab:modern_comparison} consolidates the direct and conceptual
neighbors. The relevant gap is an executable plaintext binary-image PEAK/ZERO
method in which the active reversible mapping configuration is selected by
usable mapping capacity after the exact serialized description cost needed for
deterministic recovery. Within the reviewed corpus, this study contributes the
joint treatment of serialization-aware activation, an enumerative
PEAK/flip-position wire representation, and a common executable comparison
against representative binary-image baselines.
\RevisionAddEnd

\begin{table*}[t]
\centering
\caption{Comparative overview of binary RDH and related systems}
\label{tab:modern_comparison}
\resizebox{\textwidth}{!}{
\begin{tabular}{@{}lccccccc@{}}
\toprule
\textbf{Method (Year)} & \textbf{EC (bits)} & \textbf{Net accounting} &
\textbf{PSNR (dB)} & \textbf{DRD} & \textbf{Adaptive} &
\textbf{Uniform blocks} & \textbf{Steg. eval.} \\
\midrule
Ho et al. 2009 \cite{Ho2009PatternSubstitution} & Published & Open & -- & -- & Yes & Pattern-dependent & Not extracted \\
Zhang et al. 2012 \cite{Zhang2012OptimalCodes} & Code-dependent & Reconstruction info & -- & Distortion bound & Yes & Binary cover sequence & Not extracted \\
Yin et al. 2020 \cite{Yin2020PPOCP} & 1157 & Open & 27.17 & 0.23 & Fixed & Reserved & Not extracted \\
Dong et al. 2020 \cite{Dong2020Adaptive} & $\sim$1500 & Open & $\sim$25 & $<$0.5 & Yes & Reserved & Not extracted \\
Xiao et al. 2022 \cite{Xiao2022BinaryCovers} & Variable & Reconstruction info & High & Optimized & Yes & Pixel sequence & Not extracted \\
Li et al. 2023 \cite{Li2023MultiEmbedding} & 1000 example & Reconstruction info & 49.45 & -- & Yes & Pixel sequence & Not extracted \\
Ren et al. 2023 \cite{Ren2023Sliding} & $\sim$2000 & Open & $\sim$24 & $<$1 & Fixed & Reserved & Not extracted \\
Zhang et al. 2019 \cite{Zhang2019Magnification} & 369 & Open & 32.31 & 0.08 & Fixed & Reserved & Not extracted \\
Chhajed \& Garg 2023 \cite{Chhajed2023BlockDiagonal} & $<$2000 & Open & $\sim$25 & $<$1 & Fixed & Reserved & Not extracted \\
Yang 2023 \cite{Yang2023Large} & $>$2000 & Open & $\sim$23 & $<$1 & Partial & Reserved & Not extracted \\
Huynh \& Nguyen 2025 \cite{Huynh2025BlockMapping} & 2070 & Open & 24.13 & 0.55 & Fixed & Reserved & Not extracted \\
\textbf{Proposed} & \textbf{2874 gross} & \textbf{2245 net} & \textbf{--} & \textbf{0.76} & \textbf{Yes} & \textbf{Evaluated} & \textbf{Yes} \\
\bottomrule
\end{tabular}
}
\RevisionAddBegin
\vspace{2pt}
{\footnotesize\RaggedRight EC: capacity reported by the cited study on its stated
$512\times512$ protocol. Values provide cross-protocol context because
datasets, payload definitions, and side-information accounting differ.
``Steg. eval.'': whether a steganalysis result was extracted for this overview
table; Section~\ref{subsec:steganalysis} reports the common-protocol detector
evaluation. For the proposed method, DRD is the reported common-protocol
distortion summary; per-image PSNR and changed-pixel outputs are included in
the reproducibility artifacts.\par}
\RevisionAddEnd
\end{table*}



\section{Net-Aware Learned Distance-One Block Mapping}
\label{sec:proposedmethod}

\subsection{Overview and Design Rationale}

This section specifies the complete system evaluated in
Section~\ref{sec:experimental}. The design uses the established
$3\times3$ PEAK/ZERO block-mapping family, but the active mapping tables are
chosen by usable mapping capacity after the exact serialized mapping description is
charged. The encoder and decoder are therefore defined as the pair of maps
\begin{equation}
\mathrm{Enc}(I,M)\rightarrow(I',A_{\mathrm{wire}})
\label{eq:encoder-contract}
\end{equation}
and
\begin{equation}
\mathrm{Dec}(I',A_{\mathrm{wire}})\rightarrow(I,M).
\label{eq:decoder-contract}
\end{equation}
The auxiliary stream $A_{\mathrm{wire}}$ is stored or transmitted alongside the
marked binary image $I'$ as explicit synchronization metadata. The
unconstrained adaptive threshold is retained as an ablation that characterizes
the broader capacity--distortion landscape; the deployed codec uses globally
disjoint Hamming-distance-one substitutions. Consequently, each active block
carries one bit and changes at most one pixel.

\RevisionAddBegin
The evaluated system proceeds in four deterministic steps:
\RevisionAddEnd
\begin{itemize}
    \item learns a perceptual ordering of absent distance-one patterns;
    \item enforces a globally injective PEAK/ZERO assignment;
\RevisionAddBegin
    \item selects the frequency-ordered table prefix that maximizes map-level
    net capacity; and
\RevisionAddEnd
    \item serializes the payload length, active PEAK subset, and flip positions
    with an enumerative codec.
\end{itemize}

The embedding and extraction workflows are illustrated in
Figures~\ref{fig:embedding-flow} and~\ref{fig:extraction-flow} and formalized
in Algorithms~\ref{alg:embedding} and~\ref{alg:extraction}.

\subsection{Reversible Mapping Construction}

\subsubsection*{Block representation and pattern histogram}

Let $I$ be a binary image of size $M\times N$. Complete non-overlapping
$3\times3$ blocks are visited in raster order:
\begin{equation}
B_k=\{b_{i,j}\},\quad i,j\in\{0,1,2\},\quad
k=1,\ldots,\left\lfloor\frac{M}{3}\right\rfloor
\left\lfloor\frac{N}{3}\right\rfloor .
\label{eq:block-index}
\end{equation}
The resulting block count is the product of the two complete-block counts:
\begin{equation}
B=\left\lfloor\frac{M}{3}\right\rfloor
  \left\lfloor\frac{N}{3}\right\rfloor .
\label{eq:block-count}
\end{equation}

Each block is mapped to a decimal index:
\begin{equation}
d_k = \sum_{i=0}^{8} b_i 2^{8-i},
\label{eq:pattern-index}
\end{equation}
where $b_i \in \{0,1\}$ follows a fixed row-major ordering shared by both
embedding and extraction; the first row-major bit is the most significant bit.
Rows with indices at least $3\lfloor M/3\rfloor$ and columns with indices at
least $3\lfloor N/3\rfloor$ are copied unchanged and never participate in
capacity or extraction.

Uniform blocks ($d_k = 0$ and $d_k = 511$) are reserved for the separately
evaluated Random Forest layer in Section~\ref{subsec:rf-layer}; the core
block-mapping histogram uses $d\in\{1,\ldots,510\}$.
A histogram $H(d)$ is constructed for $d \in \{1,\ldots,510\}$.

\subsubsection*{Pattern sets}

The implementation considers every non-uniform pattern occurring at least
twice as an initial PEAK candidate. Candidates are sorted by decreasing
frequency and then by numerical index. The pattern space is divided into:
\begin{itemize}
    \item \textbf{PEAK candidates} $\mathcal{P}_0=\{d:H(d)\geq2,\,
    d\notin\{0,511\}\}$,
    \item \textbf{Unused patterns (ZEROs)} $\mathcal{Z}$: patterns such that $H(d)=0$,
\RevisionAddBegin
    \item \textbf{Residual patterns}: patterns outside the embedding tables.
\RevisionAddEnd
\end{itemize}

The final active set $\mathcal{P}\subseteq\mathcal{P}_0$ is selected by the
\RevisionAddBegin
explicit map-level net-capacity rule defined below.
\RevisionAddEnd

\subsubsection*{Pattern distance analysis}

For each frequent pattern $P_i \in \mathcal{P}$ and each unused pattern $Z_j \in \mathcal{Z}$,
the Hamming distance between their $3 \times 3$ representations is computed:
\begin{equation}
D_{i,j} = \sum_{p=1}^{9} \delta(P_i[p], Z_j[p]),
\label{eq:hamming-distance}
\end{equation}
where $\delta(a,b)=1$ if $a\neq b$ and $0$ otherwise.

The set $\{D_{i,j}\}$ characterizes the distortion cost associated with mapping $P_i$ to available unused patterns.

\subsubsection*{Adaptive candidate selection}

Candidate patterns for each frequent pattern are selected adaptively.
Let $\mu_i$ and $\sigma_i$ denote the empirical mean and standard deviation of $\{D_{i,j}\}$.
The adaptive Hamming threshold and candidate set are
\begin{equation}
H_{T,i}=\mu_i+\alpha\sigma_i
\label{eq:adaptive-threshold}
\end{equation}
and
\begin{equation}
\mathcal{Z}_i = \{ Z_j \in \mathcal{Z} \mid D_{i,j} \le H_{T,i} \},
\label{eq:adaptive-candidate-set}
\end{equation}
where $\alpha$ controls the capacity--distortion trade-off.

\RevisionAddBegin
This formulation uses empirical distance statistics rather than a parametric
model for the observed distance values.
\RevisionAddEnd

\subsubsection{Globally Injective Mapping Construction}

For the unconstrained ablation, let $k_i=|\mathcal{Z}_i|$ and define
\begin{equation}
L_i = \left\lfloor \log_2 (k_i + 1) \right\rfloor
\label{eq:bits-per-occurrence}
\end{equation}
bits per occurrence. Exactly $2^{L_i}-1$ candidates are retained:
\begin{equation}
\widetilde{\mathcal Z}_i\subseteq\mathcal Z_i,\qquad
|\widetilde{\mathcal Z}_i|=2^{L_i}-1,
\label{eq:retained-zero-set}
\end{equation}
which makes the mapping
\begin{equation}
T_i:\{0,1\}^{L_i}\longleftrightarrow
\{P_i\}\cup\widetilde{\mathcal Z}_i
\label{eq:local-bijection}
\end{equation}
a genuine bijection. After a ZERO is assigned, it is removed from the
available pool; therefore the ranges of all $T_i$ are globally disjoint.

The deployed distance-one policy restricts $L_i=1$. Symbol 0 preserves
$P_i$, while symbol 1 uses the lowest predicted-cost unassigned ZERO satisfying
$D(P_i,Z)=1$. Thus every modified block flips exactly one pixel.

\subsubsection*{Deterministic block traversal order}
\label{subsec:order}

To ensure reproducibility and unambiguous extraction, a deterministic block traversal order is enforced.

The image is traversed in raster-scan order (left-to-right, top-to-bottom) at the block level.
Each block is visited exactly once during both embedding and extraction.

During embedding, frequent patterns are processed in descending order of occurrence.
However, the spatial traversal order of blocks remains unchanged.
\RevisionAddBegin
Each block is modified at most once, and block decisions are spatially local
once the mapping tables are fixed.
\RevisionAddEnd

\subsubsection*{Multi-pattern embedding strategy}
\label{subsec:multi_pattern}

Embedding scans spatial blocks once in raster order. When a block pattern is
an active PEAK, the next bit selects either the unchanged PEAK or its assigned
ZERO. The last symbol is truncated to the declared payload length during
\RevisionAddBegin
extraction; for the deployed one-bit codec, the payload boundary is
unambiguous.
\RevisionAddEnd

The total embedding capacity is:
\begin{equation}
C_{\text{gross}} = \sum_{P_i \in \mathcal{P}} n_i L_i,
\label{eq:gross-capacity}
\end{equation}
where $n_i$ denotes the number of occurrences of $P_i$.

\subsection{Embedding Procedure}
\RevisionAddBegin
The embedding procedure operationalizes the mapping construction above as a
deterministic encoder that constructs a single marked image. It first freezes
the spatial support: complete non-overlapping $3\times3$ blocks are processed
in raster order, while border pixels are copied unchanged.
This convention is used for two reasons. First, it gives the decoder the same
block sequence with coordinate-free synchronization. Second, it prevents boundary
handling from becoming hidden auxiliary information that would distort the net
payload comparison.

The next stage builds the reversible alphabet available in the current cover.
Patterns with at least two occurrences are eligible PEAKs, absent patterns are
ZEROs, and all-white/all-black uniform patterns are kept in the reserved
uniform-block pool because changing them tends to create isolated specks or holes in
document backgrounds and strokes. For each PEAK, admissible replacements are
restricted to distance-one ZEROs. This restriction makes every embedded symbol a
single-pixel block modification, keeps the restoration rule local, and allows
the auxiliary stream to store only one base-9 flip position per active table.
At this point, the CNN ranks admissible ZEROs by predicted perceptual cost;
a global disjointness check then prevents two PEAKs from sharing the same
stego pattern.

The final selection stage is the serialization-aware part of the method. After
candidate assignments have been fixed, the encoder considers frequency-ordered
prefixes of active mappings and selects the prefix that maximizes usable mapping capacity
after charging the exact serialized description of that prefix. The message bits
are then embedded during one raster scan: a payload bit 0 leaves an active PEAK
block unchanged, while a payload bit 1 replaces it by the assigned ZERO. The
declared message length $q$ is serialized with the active table description,
giving the decoder an explicit stopping point.

